Large language models produce notoriously formulaic creative writing, defaulting to predictable arcs and flat prose.
We investigate whether the \skillmix paradigm---decomposing tasks into discrete skills and generating synthetic data from novel skill combinations---transfers from instruction-following to creative story generation.
We introduce \methodname, a pipeline that extracts structured narrative components (setting, character, conflict type, theme, emotional arc, tone) from existing stories and generates new stories conditioned on novel component combinations.
Using \gptfour on 100 \rocstories, we extract a component taxonomy and generate 50 stories per condition across three settings: component-controlled, title-only baseline, and prompted baseline.
Component-controlled generation achieves near-perfect specification adherence (4.96/5.0) but incurs a statistically significant quality penalty: overall quality drops by 0.40 points ($d{=}{-}0.80$, $p{=}0.002$) compared to unconstrained baselines, with the largest degradation in language quality ($d{=}{-}0.83$) and creativity ($d{=}{-}0.46$).
Diversity metrics show modest gains (26\% larger vocabulary), but the hypothesized quality ``arbitrage''---that explicit structural control would yield better stories despite LLMs' writing limitations---is not supported.
We identify a fundamental \emph{constraint-quality tradeoff}: models allocate generation capacity to satisfying specifications at the expense of expressive prose.
Our results suggest that realizing the full potential of component-based story generation requires a two-stage approach---structural specification followed by unconstrained refinement---rather than single-pass constrained generation.
